\chapter{4位数据比较器}

\section{第一步：解决什么问题}

比较两个二进制数的大小关系。

这是数字系统中非常基础的操作：CPU的跳转指令（BEQ, BNE, BLT, BGT）本质上就是在做数值比较。

\begin{keypoint}[比较器的本质]
比较器 = 判断 $A$ 和 $B$ 的大小关系

三个输出：$A>B$、$A=B$、$A<B$（互斥，任意时刻恰好一个为1）
\end{keypoint}

\section{第二步：输入输出关系}

4位比较器：

\begin{itemize}
  \item \textbf{数据输入}：$A_3 A_2 A_1 A_0$（4位），$B_3 B_2 B_1 B_0$（4位）
  \item \textbf{级联输入}：$I_{A>B}, I_{A=B}, I_{A<B}$（用于级联多片比较器）
  \item \textbf{输出}：$Y_{A>B}, Y_{A=B}, Y_{A<B}$
\end{itemize}

\section{第三步：比较逻辑}

比较两个多位二进制数的方法：\textbf{从高位到低位逐位比较}。

\begin{enumerate}
  \item 先看最高位（MSB）$A_3$ vs $B_3$
  \begin{itemize}
    \item 如果 $A_3 > B_3$（即 $A_3=1, B_3=0$），则 $A > B$，\textbf{比较结束}
    \item 如果 $A_3 < B_3$（即 $A_3=0, B_3=1$），则 $A < B$，\textbf{比较结束}
    \item 如果 $A_3 = B_3$，则 \textbf{继续看下一位}
  \end{itemize}
  \item 然后看 $A_2$ vs $B_2$（同上逻辑）
  \item 然后看 $A_1$ vs $B_1$
  \item 最后看 $A_0$ vs $B_0$
\end{enumerate}

\textbf{这完全类比人类比较两个数字的方法}：先看最高位，相等再看次高位。

\section{第四步：逻辑推导}

一位比较器的基本单元：

\begin{itemize}
  \item $A_i > B_i$：$A_i \cdot \overline{B_i}$
  \item $A_i < B_i$：$\overline{A_i} \cdot B_i$
  \item $A_i = B_i$：$\overline{A_i \oplus B_i}$ （XNOR）
\end{itemize}

多位比较器 = 从高位到低位的级联比较：

\begin{align}
  A>B &= (A_3 > B_3) + (A_3 = B_3) \cdot (A_2 > B_2) \\
      &\quad + (A_3 = B_3) \cdot (A_2 = B_2) \cdot (A_1 > B_1) \nonumber \\
      &\quad + (A_3 = B_3) \cdot (A_2 = B_2) \cdot (A_1 = B_1) \cdot (A_0 > B_0) \nonumber
\end{align}

即：在更高位全部相等的前提下，看当前位谁大谁小。

\begin{align}
  A=B &= (A_3 = B_3) \cdot (A_2 = B_2) \cdot (A_1 = B_1) \cdot (A_0 = B_0) \\
  A<B &= \text{（对称逻辑）}
\end{align}

\section{第五步：结构化设计}

比较器内部结构 = 逐位比较单元 + 优先级逻辑：

\begin{center}
\begin{tikzpicture}
  % MSB comparison
  \node[draw, minimum width=2cm, minimum height=1.5cm] (cmp3) at (0,3) {bit3比较};
  \node[draw, minimum width=2cm, minimum height=1.5cm] (cmp2) at (0,2) {bit2比较};
  \node[draw, minimum width=2cm, minimum height=1.5cm] (cmp1) at (0,1) {bit1比较};
  \node[draw, minimum width=2cm, minimum height=1.5cm] (cmp0) at (0,0) {bit0比较};

  % Priority chain
  \draw[->] (cmp3.east) -- ++(0.5,0) |- (2,2.5) node[right] {若相等则};
  \draw[->] (cmp2.east) -- ++(0.5,0) |- (2,1.5) node[right] {若相等则};
  \draw[->] (cmp1.east) -- ++(0.5,0) |- (2,0.5) node[right] {若相等则};
\end{tikzpicture}
\end{center}

\textbf{关键设计思想}：优先级从高到低。高位比较的结果优先级最高。

\section{第六步：为什么这样设计}

\begin{enumerate}
  \item \textbf{为什么从高位到低位？}

  因为高位的权重更大。$A_3=1$表示$A \geq 8$，而$A_0=1$只表示$A \geq 1$。高位1比特的差异可以压倒低位所有比特的差异。

  这和十进制完全一样：比较 9321 和 9211，百位2>1就确定了大小，不需要再看十位和个位。

  \item \textbf{为什么需要级联输入端？}

  单芯片4位比较器不够用时（例如需要比较8位或16位数），多片级联。级联输入接收低位比较器的结果，当高位相等时输出级联输入的值。

  \item \textbf{比较器的根本限制？}

  组合逻辑比较器随着位宽增加，门延迟增加（每级AND串联）。对于大位宽（如32位、64位），现代设计中常用减法器+标志位来判断大小（但这是另一个话题了）。
\end{enumerate}

\section{第七步：考试常见变式}

\begin{enumerate}
  \item \textbf{变式1：8位比较器（用两片4位比较器级联）}

  低位比较器的输出接到高位比较器的级联输入端。当高位相等时，结果由低位决定。

  \item \textbf{变式2：只输出"相等"的比较器}

  只需要 $A=B$ 信号。实现最简单：每一位都相等（XNOR），然后全部AND。

  \item \textbf{变式3：比较器 + MUX 实现取最大值/最小值}

  比较器判断大小，MUX选择较大的（或较小的）数输出。这在排序、中值滤波等应用中很常见。

  \item \textbf{变式4：带使能的比较器}

  使能端控制比较器是否工作；不工作时所有输出为0。
\end{enumerate}

\section{第八步：VHDL实现}

\begin{lstlisting}[caption={4位比较器VHDL实现}]
library ieee;
use ieee.std_logic_1164.all;

entity comparator_4bit is
  port (
    A    : in  std_logic_vector(3 downto 0);
    B    : in  std_logic_vector(3 downto 0);
    A_gt : out std_logic;
    A_eq : out std_logic;
    A_lt : out std_logic
  );
end comparator_4bit;

architecture behavioral of comparator_4bit is
begin
  process(A, B)
  begin
    -- 默认值
    A_gt <= '0';
    A_eq <= '0';
    A_lt <= '0';

    if A > B then
      A_gt <= '1';
    elsif A = B then
      A_eq <= '1';
    else
      A_lt <= '1';
    end if;
  end process;
end behavioral;
\end{lstlisting}

\textbf{VHDL的关键优势}：直接用 \texttt{if A > B} 比较运算符，综合工具会自动生成合适的比较器硬件。不需要手写门级逻辑。

\section{第九步：Testbench验证}

\begin{lstlisting}[caption={比较器Testbench}]
library ieee;
use ieee.std_logic_1164.all;

entity comparator_4bit_tb is
end comparator_4bit_tb;

architecture test of comparator_4bit_tb is
  signal A, B : std_logic_vector(3 downto 0);
  signal A_gt, A_eq, A_lt : std_logic;

  component comparator_4bit is
    port (
      A    : in  std_logic_vector(3 downto 0);
      B    : in  std_logic_vector(3 downto 0);
      A_gt : out std_logic;
      A_eq : out std_logic;
      A_lt : out std_logic
    );
  end component;
begin
  uut: comparator_4bit port map (A, B, A_gt, A_eq, A_lt);

  stimulus: process
  begin
    -- A > B
    A <= "1010"; B <= "0101"; wait for 10 ns;
    -- A = B
    A <= "1100"; B <= "1100"; wait for 10 ns;
    -- A < B
    A <= "0011"; B <= "1100"; wait for 10 ns;
    wait;
  end process;
end test;
\end{lstlisting}

\section{第十步：如何快速解题}

\begin{examtip}[比较器快速解题法]
\begin{enumerate}
  \item \textbf{VHDL解法}：直接用 \texttt{>}, \texttt{=}, \texttt{<} 运算符，综合工具自动处理
  \item \textbf{门级解法}：从高位到低位，用"前级相等 AND 本级比较"的级联结构
  \item \textbf{级联题}：低位输出 $\to$ 高位级联输入
  \item \textbf{常考}：比较器 + MUX 组合题（找最大/最小值）
\end{enumerate}
\end{examtip}

\section{变式详解}
\chapter{4位比较器——4大变式详解}

\section{变式1：8位比较器（两片4位级联）}

\begin{detailbox}[完整教学]
\textbf{【级联方式】}
高位4位→主比较器。低位4位→从比较器，输出接主比较器的级联输入端。

\textbf{【级联逻辑】}
\begin{itemize}
  \item 主比较器比较A[7:4]和B[7:4]
  \item 如果高位相等，则看级联输入（低位比较器的结果）
  \item 低位比较器比较A[3:0]和B[3:0]，输出接到主比较器的I\_A>B, I\_A=B, I\_A<B
\end{itemize}

\textbf{【VHDL——结构级】}
\begin{lstlisting}[caption={8位比较器结构级实现}]
library ieee;
use ieee.std_logic_1164.all;

entity comparator_8bit is
  port (
    A    : in  std_logic_vector(7 downto 0);
    B    : in  std_logic_vector(7 downto 0);
    A_gt : out std_logic;
    A_eq : out std_logic;
    A_lt : out std_logic
  );
end comparator_8bit;

architecture structural of comparator_8bit is
  signal gt_lo, eq_lo, lt_lo : std_logic;

  component comparator_4bit_cascade is
    port (
      A    : in  std_logic_vector(3 downto 0);
      B    : in  std_logic_vector(3 downto 0);
      I_gt : in  std_logic;
      I_eq : in  std_logic;
      I_lt : in  std_logic;
      A_gt : out std_logic;
      A_eq : out std_logic;
      A_lt : out std_logic
    );
  end component;
begin
  -- 低位比较器
  cmp_lo: comparator_4bit_cascade
    port map(A=>A(3 downto 0), B=>B(3 downto 0),
             I_gt=>'0', I_eq=>'1', I_lt=>'0',  -- 最低位：假设之前相等
             A_gt=>gt_lo, A_eq=>eq_lo, A_lt=>lt_lo);
  -- 高位比较器
  cmp_hi: comparator_4bit_cascade
    port map(A=>A(7 downto 4), B=>B(7 downto 4),
             I_gt=>gt_lo, I_eq=>eq_lo, I_lt=>lt_lo,  -- 级联!
             A_gt=>A_gt, A_eq=>A_eq, A_lt=>A_lt);
end structural;
\end{lstlisting}

\textbf{【行为级——直接用运算符（更简洁）】}
\begin{lstlisting}[caption={8位比较器行为级实现}]
library ieee;
use ieee.std_logic_1164.all;
use ieee.std_logic_unsigned.all;

entity comparator_8bit is
  port (
    A    : in  std_logic_vector(7 downto 0);
    B    : in  std_logic_vector(7 downto 0);
    A_gt : out std_logic;
    A_eq : out std_logic;
    A_lt : out std_logic
  );
end comparator_8bit;

architecture behavioral of comparator_8bit is
begin
  A_gt <= '1' when A > B else '0';
  A_eq <= '1' when A = B else '0';
  A_lt <= '1' when A < B else '0';
end behavioral;
\end{lstlisting}
综合器自动处理任意位宽的比较！
\end{detailbox}


\section{变式2：只输出"相等"的比较器}

\begin{detailbox}[完整教学]
\textbf{【简化】}只需要$A=B$信号。每一位都相等→XNOR→全部AND。

\textbf{【VHDL】}
\begin{lstlisting}[caption={相等比较器完整实现}]
library ieee;
use ieee.std_logic_1164.all;

entity comparator_eq_only is
  port (
    A    : in  std_logic_vector(3 downto 0);
    B    : in  std_logic_vector(3 downto 0);
    A_eq : out std_logic
  );
end comparator_eq_only;

architecture behavioral of comparator_eq_only is
begin
  A_eq <= '1' when A = B else '0';
  -- 或门级：
  -- A_eq <= (A(3) xnor B(3)) and (A(2) xnor B(2)) and
  --         (A(1) xnor B(1)) and (A(0) xnor B(0));
end behavioral;
\end{lstlisting}

\textbf{【应用】}地址译码、标签匹配、Cache Tag比较。
\end{detailbox}


\section{变式3：比较器+MUX实现取最大值}

\begin{detailbox}[完整教学]
\textbf{【电路】}比较器判断A>B→输出作为MUX选择信号。A>B时MUX选A，否则选B。

\textbf{【VHDL】}
\begin{lstlisting}[caption={最大值/最小值选择器}]
library ieee;
use ieee.std_logic_1164.all;

entity max_min_selector is
  port (
    A       : in  std_logic_vector(3 downto 0);
    B       : in  std_logic_vector(3 downto 0);
    max_val : out std_logic_vector(3 downto 0);
    min_val : out std_logic_vector(3 downto 0)
  );
end max_min_selector;

architecture behavioral of max_min_selector is
begin
  max_val <= A when A > B else B;
  min_val <= A when A < B else B;
end behavioral;
\end{lstlisting}

\textbf{【综合结果】}比较器+MUX→单行代码综合成两者组合。
\end{detailbox}


\section{变式4：带使能的比较器}

\begin{detailbox}[完整教学]
\textbf{【功能】}enable=1时比较器工作；enable=0时所有输出=0。

\textbf{【VHDL】}
\begin{lstlisting}[caption={带使能的比较器}]
library ieee;
use ieee.std_logic_1164.all;

entity comparator_gated is
  port (
    A    : in  std_logic_vector(3 downto 0);
    B    : in  std_logic_vector(3 downto 0);
    en   : in  std_logic;
    A_gt : out std_logic;
    A_eq : out std_logic;
    A_lt : out std_logic
  );
end comparator_gated;

architecture behavioral of comparator_gated is
begin
  process(A, B, en)
  begin
    if en = '0' then
      A_gt <= '0'; A_eq <= '0'; A_lt <= '0';
    else
      if A > B then A_gt <= '1'; A_eq <= '0'; A_lt <= '0';
      elsif A = B then A_gt <= '0'; A_eq <= '1'; A_lt <= '0';
      else A_gt <= '0'; A_eq <= '0'; A_lt <= '1';
      end if;
    end if;
  end process;
end behavioral;
\end{lstlisting}
\end{detailbox}
